%%%%%%%%%%%%%%%%%%%%%%%%%%%%%%%%%%%%%%%%%%%%%%%%%%%%%%%%%%%%%%%%%%%%%%%%%%%%%%%
% Introduction
%%%%%%%%%%%%%%%%%%%%%%%%%%%%%%%%%%%%%%%%%%%%%%%%%%%%%%%%%%%%%%%%%%%%%%%%%%%%%%%

\subsection{From Log Loss to KL Divergence}
\label{sec:intro-logloss}

\begin{figure}[h]
    \centering
    \includegraphics[width=0.33\textwidth]{classification.png}
    \caption{A probabilistic classifier maps an input $x$ to a probability distribution over labels.}
    \label{fig:intro-classification}
\end{figure}

Information theory enters modern machine learning largely through \emph{logarithms of probabilities}.
This subsection explains why the standard classification objective
\( -\log q_\theta(y\mid x) \) is natural, and how it is exactly a KL divergence.

\paragraph{Setup.}
Let $X\in\cX$ denote an input (e.g., $\cX=\fR^d$) and $Y\in[M]=\{1,2,\dots,M\}$ a label.
A probabilistic classifier is a map
\[
    f_\theta:\cX\to\cM,
\]
where $\cM$ is the set of probability mass functions on $[M]$.
We write $f_\theta(x)=(f_\theta(x)_1,\dots,f_\theta(x)_M)$ and interpret
\(f_\theta(x)_i\approx \Pr{Y=i\mid X=x}\).
Equivalently, define the model conditional distribution
\(q_\theta(y\mid x)=f_\theta(x)_y\).

\paragraph{Negative log-likelihood (log loss).}
For one labeled example $(x,y)$, the standard loss is
\begin{equation}
    \ell_\theta(x,y) \coloneqq -\log q_\theta(y\mid x)= -\log f_\theta(x)_y.
\end{equation}
This choice is not arbitrary:

\begin{itemize}[leftmargin=2em]
    \item \textbf{Maximum likelihood.}
    Minimizing the empirical average of $-\log q_\theta(y\mid x)$ is exactly maximum likelihood estimation for a conditional model.

    \item \textbf{A coding interpretation.}
    If we were to encode the label $Y$ using a code designed for the distribution $q_\theta(\cdot\mid x)$, the ideal (Shannon) code length for outcome $y$ is $-\log q_\theta(y\mid x)$ bits.
    Thus training with log loss is equivalent to learning a model that assigns short codewords to the true labels.

    \item \textbf{Additivity over independent samples.}
    For i.i.d. data, likelihoods multiply while log-likelihoods add:
    \( -\log \prod_i q_\theta(y_i\mid x_i) = \sum_i -\log q_\theta(y_i\mid x_i)\).
\end{itemize}

\paragraph{KL divergence viewpoint.}
Represent the true label $y$ as a one-hot distribution $\delta_y\in\cM$:
\[
\delta_y(i)=\begin{cases}
1 & i=y,\\
0 & i\neq y.
\end{cases}
\]
The KL divergence between two pmfs $P,Q\in\cM$ is
\[
\KL{P}{Q}=\sum_{i:P(i)>0} P(i)\log\frac{P(i)}{Q(i)}.
\]
A direct computation gives
\begin{equation}
    \KL{\delta_y}{f_\theta(x)}
    =\log\frac{1}{f_\theta(x)_y}
    =\ell_\theta(x,y).
\end{equation}
So the log loss is exactly the KL divergence from the empirical (one-hot) label distribution to the model prediction.

However, the concern remains.
The KL divergence is also weird looking.
Why are we measuring ``distance'' between distributions in such way?
Is there another way of measuring ``distance''?
Why is it called divergence? instead of ``distance''?

To answer these questions, we need to understand the basic information measures in information theory
from entropy, mutual information, to KL divergence (or often called relative entropy).



\paragraph{Soft labels and cross entropy.}
If the target for an example is itself a distribution $p\in\cM$ (e.g., label smoothing or distillation), the natural generalization is the \emph{cross entropy}
\begin{equation}
    H(p,f_\theta(x)) \coloneqq \sum_{i=1}^M p(i)\log\frac{1}{f_\theta(x)_i}.
\end{equation}
We will prove in Section~\ref{sec:entropy-cross-entropy} that
\begin{equation}
    H(p,f_\theta(x)) = H(p) + \KL{p}{f_\theta(x)}.
\end{equation}
Since $H(p)$ is fixed by the data, minimizing cross entropy is equivalent to minimizing a KL divergence.

\subsection{Notation and Conventions}
\label{subsec:notation}

\begin{itemize}[leftmargin=2em]
    \item We use uppercase letters (e.g., $X$) for random variables and lowercase letters (e.g., $x$) for realizations.

    \item For a discrete $X$, the pmf is $p_X(x)=\Pr{X=x}$. For a continuous $X$, the pdf is $f_X(x)$.

    \item Unless stated otherwise, $\log$ is base $2$ (information measured in \emph{bits}). We write $\ln$ for the natural logarithm.

    \item Expectations and probabilities are written as $\E{\cdot}$ and $\Pr{\cdot}$.
\end{itemize}
